\documentclass[11pt]{article}
\usepackage[margin=1in]{geometry}
\usepackage{amsmath,amssymb,amsthm}
\usepackage{enumitem}
\usepackage{booktabs}
\usepackage{lmodern}
\usepackage{microtype}
\usepackage{hyperref}
\hypersetup{colorlinks=true,urlcolor=blue,linkcolor=black,citecolor=black}

\setlength{\parskip}{4pt plus 1pt}
\setlength{\parindent}{0pt}

\newtheorem{thm}{Theorem}
\newtheorem{lem}[thm]{Lemma}
\newtheorem{prop}[thm]{Proposition}
\newtheorem{cor}[thm]{Corollary}
\newtheorem{conj}[thm]{Conjecture}
\theoremstyle{definition}
\newtheorem{defn}[thm]{Definition}
\newtheorem{rem}[thm]{Remark}

\title{Locality at $n=9$ from the cyclic $1$-zeros\\
       \large via the block-rule Laurent cascade\\
       \small (satellite for the $\mathrm{Tr}(\phi^3)$ locality program)}
\author{Jonathan Valenzuela}
\date{}

\begin{document}

\maketitle

\begin{abstract}
We prove locality at $n=9$ from the cyclic $1$-zero conditions alone:
every non-local coefficient of the most general rational ansatz with
mass dimension $-2(n{-}3)$ and at most simple planar poles vanishes
once the $9$ cyclic hidden-zero constraints are imposed. The proof
combines Step-$1$ (Layer-$0$) elimination with a depth-$1$ Laurent
cascade. Step-$1$ kills $904{,}752$ of the $905{,}763$
non-triangulation size-$6$ multisets directly. The remaining
$1{,}011$ multisets fall into $113$ cyclic-orbit equivalence classes,
each non-local. Of these, $23$ orbits admit independent
single-orbit depth-$1$ cascade recipes; the other $90$ orbits form a
coupled cluster $\mathcal C_{90}$ whose depth-$1$ fingerprint matrix
has $506$ rows and column rank $90$ on cluster columns, forcing the
cluster's $90$ coefficients to vanish jointly. A complete per-orbit
audit confirming all $113$ kill mechanisms is recorded in
\cite{n9status}. \emph{This is the first polygon size at which the
cluster mechanism is required:} at $n=5,6$ Step-$1$ alone closes
locality; at $n=7,8$ every Step-$1$ survivor admits a single-orbit
cascade (cluster sizes $1$); at $n=9$ a single $90$-orbit
non-trivial cluster appears and is killed jointly by the depth-$1$
fingerprint system.
Unitarity (equality of triangulation coefficients to a single
common scalar) is independent of this paper and follows from the
$d$-subset / homogeneity-rescaling argument of
\cite{notes11,notes18,subsetvan}; locality is the load-bearing half
of the $\mathrm{Tr}(\phi^3)$ uniqueness theorem and is what this
paper closes at $n=9$.
\end{abstract}

%--------------------------------------------------------------------------
\section{Background: the three layers of kill mechanism}\label{sec:layers}
%--------------------------------------------------------------------------

We adopt the conventions of the root README. The cyclic hidden-zero
zone $\mathcal Z_{r,r+2}$ has special chord $X_{r,r+2}=:X_*$, bare
chord $X_{r+1,r-1}=-X_*$, and substitutions $X_{r+1,k}=X_{r,k}-X_*$
for $k$ in cyclic range. The general ansatz is
$B=\sum_M a_M/\prod_{c\in M}X_c$ over size-$(n{-}3)$ chord multisets
$M$.

\begin{defn}[Three layers of kill mechanism]\label{def:layers}\leavevmode\\[-1.2em]
\begin{itemize}[leftmargin=*]
\item \textbf{Step-$1$ (Layer-$0$)}: at zone $\mathcal Z_{r,r+2}$,
send $X_*\to\infty$. The leading-order coefficient identity forces
$a_M=0$ when $M\in\mathcal K_{r,r+2}$, i.e.\ when (K1) $M$ contains
neither special nor bare and (K2) $M$ contains no
(companion, substitute) pair. Holds for $M$ killable at \emph{any}
cyclic zone.

\item \textbf{Step-$2$ (equate)}: keep $X_*$ in the surviving
variables and exploit $X_{r+1,r-1}=-X_*$. This produces equality
identities $a_{M\cup\{X_{r+1,r-1}\}}=-a_{M\cup\{X_*\}}$ between
coefficients differing by a bare/special swap. Step-$2$ is the
unitarity engine; locality does \emph{not} require it.

\item \textbf{Step-$3$ (Laurent cascade)}: at zone
$\mathcal Z_{r,r+2}$, expand
\[
\frac{1}{X_{r+1,k}}=\frac{1}{X_{r,k}-X_*}
=-\frac{1}{X_*}-\frac{X_{r,k}}{X_*^2}-\frac{X_{r,k}^2}{X_*^3}-\cdots
\]
and extract the coefficient of a free-variable monomial (the
\emph{fingerprint}) at order $X_*^{-(\ell_r(M)+1)}$. This is a
linear identity among coefficients $a_M$ of multisets contributing
the same fingerprint.
\end{itemize}
\end{defn}

\begin{defn}[Cyclic orbit; cousin]\label{def:orbit}
$\mathbb Z_n$ acts on size-$(n{-}3)$ multisets by vertex shift; an
\emph{orbit} is a $\mathbb Z_n$-equivalence class. Given a Step-$3$
fingerprint equation $\sum_C\alpha_C\,a_C=0$ at
$(\mathcal Z, k, \mathcal F)$, the multisets
$\{C : \alpha_C \ne 0\}$ are the \emph{cousins} of any $M$ in the
equation; their orbit-classes are the \emph{cousin orbits}.
\end{defn}

\begin{rem}[Step-$1$ exhaustivity at low $n$]\label{rem:low-n}
At $n=4,5,6$ Step-$1$ alone kills every non-triangulation
coefficient---the count of non-triangulation Step-$1$ survivors is
$0$ at each of these polygon sizes \cite{step1stats}. Cluster and
cascade machinery are required only starting at $n=7$.
\end{rem}

%--------------------------------------------------------------------------
\section{Single-orbit cascade and block-rule cascade}\label{sec:cascades}
%--------------------------------------------------------------------------

\begin{defn}[Single-orbit cascade]\label{def:single}
Orbit $\mathcal O$ admits a \emph{single-orbit depth-$1$ cascade}
recipe iff there is a triple $(\mathcal Z, k, \mathcal F)$ with
$k=\ell_{\mathcal Z}(M)+1$ for some $M\in\mathcal O$ such that:
(i) the fingerprint equation has $a_M$ with nonzero scalar; and
(ii) every cousin of $M$ in this equation is Step-$1$-killable at
some cyclic zone.
\end{defn}

When (i)+(ii) hold, substituting cousin-zeros into the fingerprint
equation collapses it to $\alpha_M\,a_M=0$, so $a_M=0$. Cyclic
equivariance rotates the recipe onto every member of $\mathcal O$.

The mechanism that closes $n=7,8$ is: \emph{every Step-$1$ survivor
orbit admits a single-orbit recipe.} The $n=9$ data falsifies this
for $4$ orbits, and a further $14$ orbits time out under a
$600$-second search cap; the $4+14$ all lie in a single coupled
cluster, and are killed jointly by the block-rule below rather than
by individual cascades.

\begin{defn}[Coupling graph; cluster]\label{def:coupling}
The \emph{coupling graph} $G_n$ has vertex set $=$ Step-$1$ survivor
orbits at $n$. An edge $\mathcal O \sim \mathcal O'$ exists iff
there is a Step-$3$ fingerprint equation with $\mathrm{orbit}(M)
=\mathcal O$ and $\mathrm{orbit}(C)=\mathcal O'$ for some non-local
cousin $C$. A \emph{cluster} is a connected component of $G_n$.
A \emph{singleton cluster} is a connected component of size $1$
(an isolated vertex, with no non-local cousins in any Step-$3$
fingerprint equation).
\end{defn}

\begin{defn}[Block-rule cascade]\label{def:block-rule}
Let $\mathcal C=\{\mathcal O_1,\dots,\mathcal O_m\}$ be a cluster.
Enumerate all admissible depth-$1$ fingerprint recipes for orbits in
$\mathcal C$ and assemble a matrix $M_{\mathrm{FP}}^{\mathcal C}$
with rows $=$ fingerprint equations and columns $=$ cluster orbits
(\emph{cluster columns}) plus external orbits (Step-$1$ survivors
not in $\mathcal C$, plus triangulation orbits appearing as cousins).
The cluster is \emph{block-rule killable at depth-$1$} iff the
cluster sub-matrix has full column rank $m=|\mathcal C|$, given
that each non-cluster non-local survivor column is set to $0$ via
its own single-orbit cascade.
\end{defn}

The single-orbit cascade is the special case of the block-rule on a
singleton cluster. Through $n=8$ every cluster is a singleton; at
$n=9$ a non-trivial $90$-orbit cluster appears alongside $23$
singletons.

\begin{prop}[$n=7$ baseline: singleton clusters]\label{prop:n7}
At $n=7$, the seven Step-$1$ survivors form a single $\mathbb Z_7$
orbit. As a connected component of $G_7$ it is a singleton
cluster: every cousin of the orbit's fingerprint equation is
Step-$1$-killable. Equivalently, $\mathcal C=\mathcal O$ has
$|\mathcal C|=1$ and the $1\times 1$ block-rule matrix has rank $1$,
forcing $a_{\mathcal O}=0$.
\end{prop}

The proof is in \verb|paper/step2 Laurent series for hard kills/cascade_n7.tex|.
The same singleton structure (clusters of size $1$) closes
$n=8$: all $13$ cyclic orbits are singleton clusters, all killed
by their own single-orbit cascades.

%--------------------------------------------------------------------------
\section{The $n=9$ result}\label{sec:n9}
%--------------------------------------------------------------------------

Step-$1$ enumeration at $n=9$ yields $1{,}011$ non-triangulation
survivors organised into $113$ cyclic orbits (one of size $3$,
$\mathbb Z_3$-stabilised; $112$ of size $9$, free orbits). Every
representative is a non-local size-$6$ multiset: $98$ are crossing
configurations and $15$ are double-pole configurations \cite{n9status}.

\subsection{Single-orbit pass}\label{subsec:single}

The script
\verb|computations/step3_laurent/cascade_n9/cascade_kill_n9.py|
attempts a single-orbit depth-$1$ recipe for each of the $113$
representatives, with a $600$-second timeout per orbit, then a
diagnostic re-run with extended budget on the timeouts. The final
disposition:

\begin{center}
\begin{tabular}{lr}
\toprule
Outcome & count \\
\midrule
single-orbit recipe found (after re-run) & $95$ \\
search timeout (recipe still undetermined) & $14$ \\
search exhausted, no recipe found &
$4$\quad $\{\mathcal O_{22},\mathcal O_{46},\mathcal O_{88},\mathcal O_{108}\}$ \\
\bottomrule
\end{tabular}
\end{center}

The $4$ exhausted-search failures are diagnosed in
\verb|.../n9/diagnostics/diagnose_orbit_22.md|; they are not
artifacts of bounded search depth or pruning. For example,
$\mathcal O_{22}$ has representative
$M_{22}=\{X_{1,3},X_{1,4},X_{1,7},X_{1,8},X_{4,6},X_{5,7}\}$ and at
every depth $1\le k\le 4$ each fingerprint equation contains at
least one cousin that is itself a Step-$1$ survivor; specifically
$\mathcal O_{22}$ couples directly to
$\mathcal O_{88},\mathcal O_{96},\mathcal O_{102}$ at the
$\mathcal Z_{3,5}$ depth-$1$ fingerprint
$X_{3,6}/(X_{1,3}X_{1,7}X_{1,8}X_{5,7})$. Single-orbit cascade
cannot kill $\mathcal O_{22}$ in isolation; the four failures
require collective treatment.

\subsection{Cluster BFS and matrix}\label{subsec:cluster}

A breadth-first traversal of the coupling graph $G_9$ starting from
$\mathcal O_{22}$ absorbs $90$ orbits (all $4$ failures and all $14$
timeouts among them). The remaining $23$ orbits are isolated nodes
of $G_9$---singleton clusters---each killed by its own verified
single-orbit recipe.

For the $90$-orbit cluster $\mathcal C_{90}$, the depth-$1$ cascade
matrix tabulates $506$ fingerprint equations against $143$ columns
($90$ cluster orbits $+$ $32$ triangulation orbits appearing as
cousins $+$ $21$ non-cluster non-local survivor orbits). Working
over $\mathbb Q$:

\begin{center}\small
\begin{tabular}{lr}
\toprule
quantity & value\\
\midrule
cluster size $|\mathcal C_{90}|$ & $90$\\
depth-$1$ fingerprint equations & $506$\\
external columns (cluster-touched triangulation $+$ non-cluster non-local) & $32+21=53$\\
$\mathrm{rank}(M_{\mathrm{FP}}^{\mathcal C_{90}})$ on cluster columns & $\boldsymbol{90}$\\
cluster residual $r=|\mathcal C_{90}|-\mathrm{rank}$ & $\boldsymbol{0}$\\
\bottomrule
\end{tabular}
\end{center}

\begin{thm}[$n=9$ block-rule kill]\label{thm:n9-blockkill}
The depth-$1$ fingerprint sub-matrix on the $90$ cluster columns of
$M_{\mathrm{FP}}^{\mathcal C_{90}}$ has full column rank $90$.
Consequently, given that the $21$ non-cluster non-local survivor
coefficients vanish (each by its own single-orbit cascade), every
one of the $90$ cluster non-local coefficients vanishes:
\[
   a_{\mathcal O_i}=0 \qquad\forall\,\mathcal O_i\in\mathcal C_{90}.
\]
\end{thm}

\begin{proof}
The rank computation is symbolic over $\mathbb Q$ via
\texttt{sympy}; the trace is in
\verb|cluster_partial.json|. With non-cluster non-local survivor
columns set to $0$, the cluster matrix gives a homogeneous linear
system on the $90$ cluster columns plus $32$ triangulation columns.
Full column rank on cluster columns means: the assignment
``triangulations equal a common scalar $c$, non-cluster non-locals
$=0$, cluster $=0$'' satisfies the system (because
$A_n^{\mathrm{tree}}$ does, putting this vector in the nullspace),
and this is the unique cluster-coefficient solution under any
consistent triangulation assignment.
\end{proof}

\subsection{Locality at $n=9$, unconditional}\label{subsec:locality}

\begin{thm}[$n=9$ locality]\label{thm:n9-locality}
Suppose $B\big|_{\mathcal Z_r}=0$ for all nine cyclic zones at
$n=9$. Then $a_M=0$ for every non-triangulation multiset $M$ of
size $6$.
\end{thm}

\begin{proof}
Partition the $905{,}763$ non-triangulation multisets at $n=9$ by
Step-$1$ killability:
\begin{enumerate}[label=(\roman*),leftmargin=2em]
\item \textbf{Step-$1$ killed.} $904{,}752$ multisets are
Step-$1$ killable at some zone, so their coefficients vanish by
Definition~\ref{def:layers}.
\item \textbf{Singleton-cluster Step-$1$ survivors.} The remaining
$1{,}011$ multisets form $113$ cyclic orbits; $23$ of these are
singleton clusters (isolated nodes of $G_9$) and admit single-orbit
depth-$1$ cascade recipes (\S\ref{subsec:single}, audited in
\cite{n9status}). Each recipe forces its orbit's coefficient to
zero.
\item \textbf{Non-trivial cluster Step-$1$ survivors.} The remaining
$90$ orbits form $\mathcal C_{90}$. By
Theorem~\ref{thm:n9-blockkill}, all $90$ cluster non-local
coefficients vanish.
\end{enumerate}
Cases (i)--(iii) exhaust all non-triangulation multisets.
\end{proof}

\begin{cor}[$n=9$ amplitude reconstruction]\label{cor:n9-amp}
Combined with the $d$-subset / homogeneity-rescaling unitarity
argument of \cite{notes11,notes18,subsetvan}, which equates all
triangulation coefficients to a common scalar $c$,
$B=c\cdot A_9^{\mathrm{tree}}$.
\end{cor}

The unitarity input is logically independent of this paper.
Locality (Theorem~\ref{thm:n9-locality}) closes the load-bearing
half of the $\mathrm{Tr}(\phi^3)$ uniqueness claim at $n=9$;
the $d$-subset argument closes the unitarity half.

\subsection{Comparison with $n=7,8$}\label{subsec:compare}

\begin{center}
\begin{tabular}{lccc}
\toprule
& $n=7$ & $n=8$ & $n=9$ \\
\midrule
Step-$1$ non-tri survivors & $7$ & $100$ & $1{,}011$ \\
\quad cyclic orbits & $1$ & $13$ & $113$ \\
singleton clusters $|\mathcal C|=1$ & $1$ & $13$ & $23$ \\
non-trivial clusters $|\mathcal C|>1$ & $0$ & $0$ & $1$ (size $90$) \\
fingerprint equations on largest cluster & $1$ & $1$ & $506$ \\
matrix rank on cluster columns & $1$ & $1$ & $90$ \\
cluster residual $r$ & $0$ & $0$ & $0$ \\
\bottomrule
\end{tabular}
\end{center}

The growth from $n=7,8$ to $n=9$ is qualitative: at $n\le 8$ the
cluster mechanism is trivial ($|\mathcal C|=1$ everywhere) and the
locality proof reduces to the singleton cascade. At $n=9$ a single
non-trivial cluster of $90$ orbits appears, and is killed jointly
by the depth-$1$ system. The $5\!\times$ ratio of fingerprint
equations to cluster orbits ($506/90$) overdetermines the cluster
and gives residual $r=0$. Whether this overdetermination persists
at higher $n$, or whether a non-trivial residual ever appears,
is the central all-$n$ open question (\S\ref{sec:remaining}).

%--------------------------------------------------------------------------
\section{Geometric interpretation}\label{sec:geom}
%--------------------------------------------------------------------------

The cascade mechanism is the algebraic shadow of a geometric
operation. At zone $\mathcal Z_{r,r+2}$, the substitution
$X_{r+1,k}=X_{r,k}-X_*$ encodes the \emph{ear-collapse} that
identifies vertex $r{+}1$ with the rest of the polygon. The bare
relation $X_{r+1,r-1}=-X_*$ is the \emph{$4$-gon kinematic
identity}: the only chord relation at $n=4$ is $X_{1,3}=-X_{2,4}$,
and every higher-$n$ collapse uses this identity at one local
cyclic position.

The Laurent expansion in $1/X_*$ records the \emph{kinematic
memory} of the collapse: each non-bare substitute $X_{r+1,k}$
remembers its companion $X_{r,k}$ at order $X_*^{-(j+1)}$
(coefficient $X_{r,k}^{j}$). Depth-$d$ probes the $d$-th layer of
this memory.

A non-local multiset killable by a depth-$d$ cascade has $d$
chord-pairs whose crossing structure interacts with the relevant
ear-collapse. A multiset \emph{not} killable by any single-zone
cascade---the cluster case---has crossings coupled across multiple
cyclic orbits: configurations reachable from each other by a
single ear-collapse move. The cluster is the connected component
of the associahedron \emph{face graph} restricted to ear-collapse
moves carried by the chosen non-local cousins.

The fact that the depth-$1$ system on a cluster has rank
$\ge|\mathcal C|-1$ is, geometrically, the statement that a
connected face of the associahedron has only finitely many
``free directions'' left after the cyclic-$1$-zero constraints are
imposed. At small $n$ the face is small enough that a $1$-form
might be needed to close the last degree of freedom (a possibility
not yet observed in the data through $n=9$); at $n=9$ the
$90$-orbit face is rigidified by sheer constraint count.

%--------------------------------------------------------------------------
\section{Refined all-$n$ locality conjecture}\label{sec:conj}
%--------------------------------------------------------------------------

The $n=9$ data sharpens the all-$n$ statement.

\begin{conj}[Block-rule cascade locality]\label{conj:blockrule}
For every $n\ge 7$ and every cluster $\mathcal C$ of Step-$1$
survivor orbits at $n$:
\begin{enumerate}[label=\textup{(\arabic*)}]
\item the cluster $\mathcal C$ is finite;
\item there exists $r=r(\mathcal C)\ge 0$ with
$\mathrm{rank}(M_{\mathrm{FP}}^{\mathcal C})=|\mathcal C|-r$ on
cluster columns;
\item if $r\ge 1$, the residual $r$-dimensional space is closed
by $r$ higher-Laurent-order anchor equations.
\end{enumerate}
Equivalently: every non-local coefficient of an $n$-point ansatz
satisfying the $n$ cyclic $1$-zero conditions is forced to vanish
by Step-$1$ together with the depth-$1$ block-rule cascade and a
finite collection of higher-Laurent anchors.
\end{conj}

\begin{cor}[All-$n$ locality, conditional]\label{cor:allnloc}
Conjecture~\ref{conj:blockrule}, together with the Step-$1$
exhaustivity recorded in \cite{step1stats} and the
triangulation-survives-Step-$1$ lemma of \cite{tri-survives},
implies $a_M=0$ for every non-triangulation $M$ at every $n$.
\end{cor}

The conjecture is verified at $n=4,5,6,7,8,9$:
\begin{itemize}[leftmargin=*,topsep=2pt,itemsep=2pt]
\item $n=4,5,6$: every non-local Step-$1$ killable; no Step-$3$ work
needed.
\item $n=7,8$: every cluster is a singleton ($|\mathcal C|=1$),
$r=0$.
\item $n=9$: $|\mathcal C_{90}|=90$, $r=0$
(Theorem~\ref{thm:n9-blockkill}); plus $23$ singleton clusters,
each $r=0$.
\end{itemize}

The pattern across these data points is consistent with the
prediction $r=0$ throughout: the fingerprint-equation count
overdetermines clusters whenever they grow beyond singletons.
A non-trivial residual ($r\ge 1$) has not been observed in any
verified case and may not appear at any $n$.

%--------------------------------------------------------------------------
\section{What remains for the all-$n$ locality theorem}\label{sec:remaining}
%--------------------------------------------------------------------------

The proof architecture is fully laid out. The remaining mathematical
obligations for the all-$n$ locality theorem are:

\begin{enumerate}[leftmargin=*,topsep=2pt,itemsep=2pt]
\item \textbf{Cluster finiteness.} Every cluster $\mathcal C$ at
every $n$ is finite. Likely follows from finiteness of
size-$(n{-}3)$ chord multisets, but should be recorded as a lemma.

\item \textbf{Cluster rank lower bound.} Show
$\mathrm{rank}(M_{\mathrm{FP}}^{\mathcal C})\ge|\mathcal C|-1$ on
cluster columns for all $n$. This is the core combinatorial-geometric
content. A generating-function argument counting fingerprint
equations vs cluster orbits is plausible; an associahedron-face
rigidity argument is the structural counterpart.

\item \textbf{Anchor existence when $r=1$.} If a cluster ever has a
$1$-dimensional residual at any $n$, exhibit a higher-Laurent-order
anchor that closes it. Through $n=9$ no such case has been observed,
but the all-$n$ proof should accommodate the possibility.
\end{enumerate}

These are the only obligations for locality at all $n$. Unitarity
(equality of triangulation coefficients) is independent and is
handled by the $d$-subset / homogeneity-rescaling argument; see
\cite{notes11,notes18,subsetvan}.

%--------------------------------------------------------------------------
\section*{Files and reproducibility}
%--------------------------------------------------------------------------

\begin{itemize}[leftmargin=*]
\item \verb|computations/step1_layer0_kill/kill_enumeration/|
--- Step-$1$ enumeration through $n=9$.
\item \verb|computations/step3_laurent/cascade_n9/cascade_kill_n9.py|
--- the single-orbit cascade engine.
\item \verb|computations/step4_laurent_block_analysis/n9/scripts/|
--- the cluster-BFS, matrix-rank, and anchor-search scripts.
\item \verb|computations/step4_laurent_block_analysis/n9/outputs/|
--- \verb|cluster_partial.json| (the rank-$90$ matrix trace),
\verb|cluster_partial.md| (human-readable summary),
\verb|results_cascade_n9_reps.json| (per-orbit cascade outcomes),
\verb|n9_locality_status.md| (\textbf{the consolidated per-orbit
locality audit}).
\item \verb|computations/step4_laurent_block_analysis/n9/diagnostics/|
--- per-anomaly walkthroughs for orbits $22$, $25$, $28$, $34$,
$35$, $46$, $88$, $108$.
\end{itemize}

\begin{thebibliography}{99}
\bibitem{n9status}
\emph{Consolidated $n=9$ locality status.} This repository,
\verb|computations/step4_laurent_block_analysis/n9/outputs/|
\verb|n9_locality_status.md|.

\bibitem{step1stats}
\emph{Step-$1$ kill statistics and asymptotic locality.} Companion
satellite, this repository,
\verb|paper/step1 Kill technique and statistics/step-one-kill-statistics/|.

\bibitem{tri-survives}
\emph{Triangulations are not Step-$1$ killable.} Companion
satellite, this repository,
\verb|paper/step1 Kill technique and statistics/step-one-doesnt-kill-triangulations/|.

\bibitem{subsetvan}
\emph{Proof of Rodina's claim $B=0\Rightarrow B_i=0$.} Companion
note, this repository,
\verb|paper/Details missing in Hidden zero -> unitarity Rodina proof/Proof of Rodina claim B=0->B_i = 0/|.

\bibitem{cascade_n7}
\emph{The Laurent-tail kill at $n=7$: seven layer-$0$ survivors
and how they die.} Companion satellite, this repository,
\verb|paper/step2 Laurent series for hard kills/|.

\bibitem{notes11}
Handwritten note \#11, \emph{finale-d-subset-induction}, this
repository, \verb|notes/|.

\bibitem{notes18}
Handwritten note \#18, \emph{d-subset rigorous proof}, this
repository, \verb|notes/|.

\bibitem{Rodina}
L.~Rodina, \emph{Hidden zeros are equivalent to enhanced
ultraviolet scaling and lead to unique amplitudes in
$\mathrm{Tr}(\phi^3)$ theory.} arXiv:2406.04234, 2024.

\bibitem{ABHY}
N.~Arkani-Hamed, Q.~Cao, J.~Dong, C.~Figueiredo, S.~He,
\emph{Hidden zeros for particle/string amplitudes and the unity of
colored scalars, pions, and gluons.} arXiv:2312.16282, 2023.
\end{thebibliography}

\end{document}
